\section{Conclusion}
\label{sec:conclusion}

EdgeFault-Bench reframes the credibility problem of cross-condition fault
diagnosis as a software-engineering problem and answers it the way software
engineering answers such questions: with a specification, an architecture
that enforces it, and evidence that the enforcement works. The experiment
lifecycle is a guarded state machine whose six integrity invariants are
model-checked clean over all reachable states, exercised by oracle-based
property tests, and attacked by a fault-injection catalogue that every
guard survives fail-closed. Heterogeneous datasets---bearings, a gearbox,
acoustic and vibration archives with formats, rates, and quirks that
include an unstated sampling rate and an upstream misspelling---integrate
as plugins while the kernel stays free of dataset logic, asserted by test.
Results carry content-hashed provenance, complete seed bookkeeping with
failures retained, and recording-level bootstrap uncertainty; reproduction
holds exactly across Python versions and is gated across operating systems
in CI. The costs are measured, and the boundaries are stated where evidence
ends: the formal layer certifies the design, third-party adoption evidence
is solicited rather than claimed, and non-protocol demonstration runs are
labelled as such. We believe the wider lesson generalizes beyond fault
diagnosis: when an empirical field's conclusions are fragile, the benchmark
infrastructure itself deserves specification, enforcement, and adversarial
verification---not just standardization of datasets and metrics.
